%!TEX root = rapport.tex

\chapter{Relevant Properties of $\beta$-reduction Graphs}

There are four main properties that need to be adressed when drawing
$\beta$-reduction graphs.

\section{Size}

Reduction graphs vary greatly in size, from one node to an infinite amount of
nodes and edges. For example, $G_\beta(\lam{x}{x})$ contains one node and zero
edges since there are no reducts in the term. On the other hand, the graph
$G_\beta((\lam{x}{xxx})(\lam{x}{xxx}))$ contains infinitely many nodes and
edges. A drawing algorithm for such graphs must therefore be able to cope with
small graphs as well as very large graphs. 


% Flere eksempler?

\section{Direction}

The graphs are directed, so arrow heads need to be drawn. Since an edge can be
drawn spanning a lot of screen area, the type of arrow has to be designed such
that it is clear to see where an edge starts from the arrow head alone and
vice versa. An alternative to arrows is drawing the edges with a colour
gradient. In this approach, an edge would be couloured e.g. green at the start
and red at the end, interpolating the colours on the edge line. The benefit of
such an approach would be that it is easy to see tendencies if there is a
visual clutter of edges having the same general direction.

\section{Planarity}

As a consequence of Lemma \ref{lem:complete} we know that in the general case,
a reduction graph can contain as its subgraph $K_5$, whence the graph is
non-planar according to Kuratowski's theorem. Thus, it is in general not
possible to planarize a reduction graph.

\begin{lemma}\label{lem:complete}
	$\forall n\in\N:\ \exists M\in\Lambda$ such that the $\beta$-reduction graph of $M$
	contains the complete (undirected) graph of $n$ nodes, $K_n$.
\end{lemma}
\begin{proof}[Sketch of proof]
	Let $n>1\in\N$. (The case $n=1$ is trivial, as $K_1$ contains no edges.)
	Let $M\in\Lambda$ be
	\begin{equation*}
		(\lam{x_1}{y}) ((\lam{x_2}{y}) ((\lam{x_3}{y})\hdots(\lam{x_{n-1}}{y})))
	\end{equation*}
	$M$ has $n-1$ distinct reducts of the form
	\begin{align*}
		&(\lam{x_1}{y}), \\
		&(\lam{x_1}{y})((\lam{x_2}{y}) y), \\
		&(\lam{x_1}{y})((\lam{x_2}{y}) ((\lam{x_3}{y}) y)),\ \hdots
	\end{align*}
	Call them $r_1,\hdots,r_{n-1}$ from smallest to greatest. For every
	$1\leq i < j \leq n-1$ we have $r_j\barrow r_i$. Thus, since $M\barrow r_i$ 
	for $1\leq i \leq n-1$, we have for every pair $(N_0,N_1)$ in the $\beta$-reduction 
	graph that $N_0\barrow N_1$ or $N_1\barrow N_0$, whence the underlying 
	undirected graph of the reduction graph is $K_n$.
\end{proof}

\section{In- and out-degree of nodes}

The in- and out-degree of nodes varies a lot; the amount of redexes in terms
can be very different from term to term, resulting in varying out-degrees,
while some nodes can have a potientially infinitely high in-degree. For
example, in the reduction graph for the term $(\lam{a}{b}) (\lam{x}{(xxx)}
\lam{y}{(yyy)})$, every node of the reduction graph has an out-degree equal to
2 and an in-degree of 1, except the one node representing the normal form ---
that node has an infinitely high in-degree and an out-degree of 0.


